% GENERATED FILE. Edit canonical lessons, then run npm run generate:books.
% canonical-source-hash: fnv1a64:fa4b4a98a5d5ba6d
% canonical-lessons: FA-C363-eddea, FA-C363-eddea-kardan, FA-C363-esbat-kardan, FA-C363-sabet-shodan, FA-C363-shavahed

\chapter{Claims and proof}
\label{ch:fa-claims-and-proof}
\hlchaptermodality{363}

\begin{chapteropening}
\textbf{By the end of this chapter:} \emph{I can make a claim, say what someone claims, prove a point and talk about the evidence.}

Complete the last lesson of chapter 363: I can make a claim, say what someone claims, prove a point and talk about the evidence.
\end{chapteropening}

\section[\texorpdfstring{\fa{ادعا} (edde'â)}{ادعا (edde'â)}]{\fa{ادعا} (edde'â) — a claim}

\label{lesson:FA-C363-eddea}



\begin{quote}
Before the new one: say the Persian for logical, then the Persian for to convince.
\end{quote}

\subsection*{You'll want to know: \fa{ادعا}}
\textbf{\fa{ادعا}} (\emph{edde'â}) — \textquotedblleft{}a claim\textquotedblright{}.

\begin{quote}
\fa{این} \fa{ادعا} \fa{درست} \fa{نیست}.
\end{quote}

\begin{quote}
\emph{in edde'â dorost nist.}
\end{quote}

\begin{quote}
This claim is not true.
\end{quote}

Plural \emph{edde'âhâ}. Formal; also a claim in court.

\begin{cousinweb}
\textbf{\fa{ادعا}} is usually traced to Arabic \emph{iddi'ā'} \textquotedblleft{}a claim, a pretension\textquotedblright{}; Persian drops the final hamza.
\end{cousinweb}

\begin{grammarlens}[title={What you've built}]
One more word for telling, explaining and describing at B1.
\end{grammarlens}

\subsection*{Your turn}
\begin{itemize}
\raggedright
  \item \emph{Say it:} \emph{edde'â}
  \item \emph{Say it:} \emph{edde'â}, once more
  \item \emph{Say it:} a sentence of your own with \emph{edde'â}
  \item \emph{Recall:} say the Persian for logical, then the Persian for to convince, then say \emph{edde'â} again
\end{itemize}

\begin{tcolorbox}[breakable,colback=teal!4,colframe=teal!35!black,title={Before you move on}]
What does \fa{ادعا} mean? (\textquotedblleft{}A claim\textquotedblright{}.) Say it once more.
\end{tcolorbox}
\section[\texorpdfstring{\fa{ادعا} \fa{کردن} (edde'â kardan)}{ادعا کردن (edde'â kardan)}]{\fa{ادعا} \fa{کردن} (edde'â kardan) — to claim}

\label{lesson:FA-C363-eddea-kardan}



\begin{quote}
Before the new one: say the Persian for to convince, then the Persian for a claim.
\end{quote}

\subsection*{You'll want to know: \fa{ادعا} \fa{کردن}}
\textbf{\fa{ادعا} \fa{کردن}} (\emph{edde'â kardan}) — \textquotedblleft{}to claim\textquotedblright{}.

\begin{quote}
\fa{او} \fa{ادعا} \fa{میکند} \fa{که} \fa{آن} \fa{مرد} \fa{را} \fa{دیده} \fa{است}.
\end{quote}

\begin{quote}
\emph{u edde'â mikonad ke ân mard râ dide ast.}
\end{quote}

\begin{quote}
He claims that he has seen that man.
\end{quote}

\emph{Edde'â} \textquotedblleft{}a claim\textquotedblright{} + \emph{kardan} (stem \emph{kon-}), followed by \emph{ke} and a clause.

\begin{cousinweb}
\textbf{\fa{ادعا} \fa{کردن}} joins \emph{edde'â}, usually traced to Arabic \emph{iddi'ā'} \textquotedblleft{}a claim\textquotedblright{}, to the inherited Persian verb \emph{kardan} \textquotedblleft{}to do\textquotedblright{}.
\end{cousinweb}

\begin{grammarlens}[title={What you've built}]
One more word for telling, explaining and describing at B1.
\end{grammarlens}

\subsection*{Your turn}
\begin{itemize}
\raggedright
  \item \emph{Say it:} \emph{edde'â kardan}
  \item \emph{Say it:} \emph{edde'â kardan}, once more
  \item \emph{Say it:} a sentence of your own with \emph{edde'â kardan}
  \item \emph{Recall:} say the Persian for to convince, then the Persian for a claim, then say \emph{edde'â kardan} again
\end{itemize}

\begin{tcolorbox}[breakable,colback=teal!4,colframe=teal!35!black,title={Before you move on}]
What does \fa{ادعا} \fa{کردن} mean? (\textquotedblleft{}To claim\textquotedblright{}.) Say it once more.
\end{tcolorbox}
\section[\texorpdfstring{\fa{اثبات} \fa{کردن} (esbât kardan)}{اثبات کردن (esbât kardan)}]{\fa{اثبات} \fa{کردن} (esbât kardan) — to prove}

\label{lesson:FA-C363-esbat-kardan}



\begin{quote}
Before the new one: say the Persian for a claim, then the Persian for to claim.
\end{quote}

\subsection*{You'll want to know: \fa{اثبات} \fa{کردن}}
\textbf{\fa{اثبات} \fa{کردن}} (\emph{esbât kardan}) — \textquotedblleft{}to prove\textquotedblright{}.

\begin{quote}
\fa{او} \fa{نمیتواند} \fa{حرفش} \fa{را} \fa{اثبات} \fa{کند}.
\end{quote}

\begin{quote}
\emph{u nemitavânad harfash râ esbât konad.}
\end{quote}

\begin{quote}
He cannot prove what he says.
\end{quote}

\emph{Esbât} \textquotedblleft{}proof\textquotedblright{} + \emph{kardan}. Formal; everyday speech prefers \emph{sâbet kardan}. In Persian, \fa{ث} sounds like s.

\begin{cousinweb}
\textbf{\fa{اثبات}} comes from Arabic \emph{ithbāt} \textquotedblleft{}proving, establishing\textquotedblright{} (EtymDB); Persian says the Arabic \emph{th} as \emph{s}. With \emph{kardan} \textquotedblleft{}to do\textquotedblright{} it becomes a verb.
\end{cousinweb}

\begin{grammarlens}[title={What you've built}]
One more word for telling, explaining and describing at B1.
\end{grammarlens}

\subsection*{Your turn}
\begin{itemize}
\raggedright
  \item \emph{Say it:} \emph{esbât kardan}
  \item \emph{Say it:} \emph{esbât kardan}, once more
  \item \emph{Say it:} a sentence of your own with \emph{esbât kardan}
  \item \emph{Recall:} say the Persian for a claim, then the Persian for to claim, then say \emph{esbât kardan} again
\end{itemize}

\begin{tcolorbox}[breakable,colback=teal!4,colframe=teal!35!black,title={Before you move on}]
What does \fa{اثبات} \fa{کردن} mean? (\textquotedblleft{}To prove\textquotedblright{}.) Say it once more.
\end{tcolorbox}
\section[\texorpdfstring{\fa{ثابت} \fa{شدن} (sâbet shodan)}{ثابت شدن (sâbet shodan)}]{\fa{ثابت} \fa{شدن} (sâbet shodan) — to be proven}

\label{lesson:FA-C363-sabet-shodan}



\begin{quote}
Before the new one: say the Persian for to claim, then the Persian for to prove.
\end{quote}

\subsection*{You'll want to know: \fa{ثابت} \fa{شدن}}
\textbf{\fa{ثابت} \fa{شدن}} (\emph{sâbet shodan}) — \textquotedblleft{}to be proven\textquotedblright{}.

\begin{quote}
\fa{حرف} \fa{او} \fa{ثابت} \fa{شد}.
\end{quote}

\begin{quote}
\emph{harf-e u sâbet shod.}
\end{quote}

\begin{quote}
What she said was proven.
\end{quote}

\emph{Sâbet} \textquotedblleft{}fixed, firm\textquotedblright{} + \emph{shodan} \textquotedblleft{}to become\textquotedblright{} (present stem \emph{shav-}). \emph{Sâbet kardan} — to prove.

\begin{cousinweb}
\textbf{\fa{ثابت}} comes from Arabic \emph{thābit} \textquotedblleft{}firm, constant, unchanging\textquotedblright{} (EtymDB), from the same Arabic family as \emph{ithbāt}. A claim that is proven has become firm.
\end{cousinweb}

\begin{grammarlens}[title={What you've built}]
One more word for telling, explaining and describing at B1.
\end{grammarlens}

\subsection*{Your turn}
\begin{itemize}
\raggedright
  \item \emph{Say it:} \emph{sâbet shodan}
  \item \emph{Say it:} \emph{sâbet shodan}, once more
  \item \emph{Say it:} a sentence of your own with \emph{sâbet shodan}
  \item \emph{Recall:} say the Persian for to claim, then the Persian for to prove, then say \emph{sâbet shodan} again
\end{itemize}

\begin{tcolorbox}[breakable,colback=teal!4,colframe=teal!35!black,title={Before you move on}]
What does \fa{ثابت} \fa{شدن} mean? (\textquotedblleft{}To be proven\textquotedblright{}.) Say it once more.
\end{tcolorbox}
\section[\texorpdfstring{\fa{شواهد} (shavâhed)}{شواهد (shavâhed)}]{\fa{شواهد} (shavâhed) — evidence}

\label{lesson:FA-C363-shavahed}



\begin{quote}
Before the new one: say the Persian for to prove, then the Persian for to be proven.
\end{quote}

\subsection*{You'll want to know: \fa{شواهد}}
\textbf{\fa{شواهد}} (\emph{shavâhed}) — \textquotedblleft{}evidence\textquotedblright{}.

\begin{quote}
\fa{شواهد} \fa{زیادی} \fa{وجود} \fa{دارد}.
\end{quote}

\begin{quote}
\emph{shavâhed-e ziyâdi vojud dârad.}
\end{quote}

\begin{quote}
There is a lot of evidence.
\end{quote}

An Arabic broken plural of \emph{shâhed} \textquotedblleft{}a witness; a piece of evidence\textquotedblright{}.

\begin{cousinweb}
\textbf{\fa{شواهد}} is usually traced to Arabic \emph{shawāhid}, a broken plural of \emph{shāhid} \textquotedblleft{}a witness\textquotedblright{} (EtymDB lists \emph{shāhid} as the source of Persian \emph{shâhed}). Evidence is what bears witness.
\end{cousinweb}

\begin{grammarlens}[title={What you've built}]
One more word for telling, explaining and describing at B1.
\end{grammarlens}

\subsection*{Your turn}
\begin{itemize}
\raggedright
  \item \emph{Say it:} \emph{shavâhed}
  \item \emph{Say it:} \emph{shavâhed}, once more
  \item \emph{Say it:} a sentence of your own with \emph{shavâhed}
  \item \emph{Recall:} say the Persian for to prove, then the Persian for to be proven, then say \emph{shavâhed} again
\end{itemize}

\begin{tcolorbox}[breakable,colback=teal!4,colframe=teal!35!black,title={Before you move on}]
What does \fa{شواهد} mean? (\textquotedblleft{}Evidence\textquotedblright{}.) Say it once more.
\end{tcolorbox}
